\documentclass[12pt,a4paper]{article}
% Drake_preamble.tex - LaTeX Preamble Template
% 作者: 謝慕揚, MD, PhD, FESC
% 
% 使用說明:
% 1. 表格請使用 longtable，不要有垂直分隔線 |
% 2. 表格內換行使用 \newline，不要使用 \\
% 3. 藥物名稱、檢驗、檢查請維持英文
% 4. Reference 格式請使用 \href 加上驗證過的 DOI

% Including essential packages for Chinese typesetting
\usepackage{xeCJK}
\usepackage{fontspec}
% Configuring Chinese font with Noto Serif CJK TC, fallback to SimSun
\IfFontExistsTF{Noto Serif CJK TC}{
  \setCJKmainfont{Noto Serif CJK TC}
  \setCJKsansfont{Noto Serif CJK TC}
  \setCJKmonofont{Noto Serif CJK TC}
}{\setCJKmainfont{SimSun}
  \setCJKsansfont{SimSun}
  \setCJKmonofont{SimSun}
}

% Including other necessary packages
\usepackage{geometry}
\usepackage{titlesec}
\usepackage{enumitem}
\usepackage{xcolor}
\usepackage{colortbl}
\usepackage{tcolorbox}
\usepackage{booktabs}
\usepackage{longtable}
\usepackage{array}
\usepackage{graphicx}
\usepackage{tikz}
\usepackage{fancyhdr}
\usepackage{multicol}
\usepackage{datetime2}
\usepackage{hyperref}
\usepackage{mdframed}
\usepackage{amsmath}
\usepackage{amssymb}
\usepackage{multirow}

\hypersetup{
    colorlinks=true,
    linkcolor=emergency_blue,
    citecolor=emergency_blue,
    urlcolor=emergency_blue,
    pdfborder={0 0 0}
}

% Defining page geometry
\geometry{
  top=2.5cm,
  bottom=2.5cm,
  left=2cm,
  right=2cm,
  headheight=25pt
}

% Adjusting topmargin to compensate for increased headheight
\addtolength{\topmargin}{-10pt}

% Defining colors
\definecolor{emergency_red}{RGB}{186,24,27}
\definecolor{emergency_blue}{RGB}{0,114,188}
\definecolor{emergency_gray}{RGB}{120,120,120}
\definecolor{highlight_yellow}{RGB}{255,250,205}

% Setting title formats
\titleformat{\section}[block]{\Large\bfseries\color{emergency_red}}{\thesection}{10pt}{}
\titleformat{\subsection}[block]{\large\bfseries\color{emergency_blue}}{\thesubsection}{8pt}{}
\titleformat{\subsubsection}[block]{\normalsize\bfseries\color{emergency_gray}}{\thesubsubsection}{6pt}{}

% Defining custom environment for important notes
\newmdenv[
    linecolor=emergency_red,
    backgroundcolor=highlight_yellow,
    linewidth=2pt,
    topline=true,
    bottomline=true,
    leftline=true,
    rightline=true,
    innertopmargin=10pt,
    innerbottommargin=10pt,
    innerleftmargin=10pt,
    innerrightmargin=10pt
]{importantbox}

% Defining note command with tcolorbox
\newcommand{\note}[1]{
    \par
    \noindent
    \begin{tcolorbox}[
        colback=emergency_blue!5!white,
        colframe=emergency_blue,
        title=\textbf{重點提醒},
        fonttitle=\bfseries,
        boxrule=0.5pt,
        arc=3pt,
        left=6pt,
        right=6pt,
        top=6pt,
        bottom=6pt
    ]
    #1
    \end{tcolorbox}
}

% Key findings box
\newcommand{\keyfinding}[1]{
    \par
    \noindent
    \begin{tcolorbox}[
        colback=yellow!10!white,
        colframe=orange!75!black,
        title=\textbf{核心發現摘要},
        fonttitle=\bfseries,
        boxrule=0.5pt,
        arc=3pt
    ]
    #1
    \end{tcolorbox}
}

% Defining custom date format for ISO 8601
\DTMnewdatestyle{iso}{
  \renewcommand{\DTMdisplaydate}[4]{\number##1-\number##2-\number##3}
  \renewcommand{\DTMDisplaydate}[4]{\number##1-\number##2-\number##3}
}
\DTMsetdatestyle{iso}
\newcommand{\mydate}{\DTMtoday}


% Custom header for this document
\pagestyle{fancy}
\fancyhf{}
\fancyhead[L]{\textcolor{emergency_red}{AJRCCM 教學講義}}
\fancyhead[R]{\textcolor{emergency_blue}{AIRFLOW-3 TLD in COPD}}
\fancyfoot[C]{\textcolor{emergency_gray}{\thepage}}
\renewcommand{\headrulewidth}{1pt}
\renewcommand{\headrule}{\hbox to\headwidth{\color{emergency_red}\leaders\hrule height \headrulewidth\hfill}}

\begin{document}

\begin{center}
{\LARGE\bfseries\textcolor{emergency_red}{American Journal of Respiratory and Critical Care Medicine}}\\[0.5cm]
{\Large\textbf{AIRFLOW-3: Randomized Sham-Controlled Trial of Targeted Lung Denervation in Patients with COPD}}\\[0.3cm]
{\large AIRFLOW-3：COPD 病人標靶肺去神經術之隨機假對照試驗}\\[0.5cm]
{\normalsize Am J Respir Crit Care Med 2025;211:2318-2329. DOI: 10.1164/rccm.202502-0404OC}\\[0.3cm]
{\normalsize 整理：謝慕揚 MD, PhD, FESC \quad 日期：\mydate}
\end{center}

\keyfinding{
\textbf{核心訊息：}AIRFLOW-3 未達到主要終點（減少 COPD exacerbation），但發現 TLD 可保留肺功能並改善呼吸困難。Post hoc 分析識別出 responder phenotype。

\textbf{主要結果：}
\begin{itemize}
    \item \textbf{Primary endpoint：}Moderate/severe exacerbation HR 1.27 (95\% CI, 0.988-1.628)，\textbf{未達統計顯著}
    \item \textbf{Dyspnea 改善：}TLD 組 TDI $\geq$1 point improvement 35.4\% vs Sham 24.1\% (P=0.021)
    \item \textbf{肺功能保留：}TLD 組 FEV$_1$ 無下降，Sham 組顯著下降 (P<0.001)
\end{itemize}
}

\tableofcontents
\newpage

%==============================================================
\section{文獻概述}
%==============================================================

本篇為 American Journal of Respiratory and Critical Care Medicine 的 Original Article，由英國 Royal Brompton Hospital 的 Shah 等人代表 AIRFLOW-3 Study Group 撰寫，報告 Targeted Lung Denervation (TLD) 在 COPD 病人的 randomized sham-controlled trial 結果。

\subsection{作者資訊}
\textbf{通訊作者：}Gerard J. Criner, MD\\
\textbf{機構：}Department of Thoracic Medicine and Surgery, Lewis Katz School of Medicine at Temple University, Philadelphia, PA, USA

\subsection{研究註冊}
ClinicalTrials.gov: NCT03639051

%==============================================================
\section{背景}
%==============================================================

\subsection{COPD 治療的未滿足需求}

\begin{itemize}
    \item 許多 COPD 病人即使接受 optimal medical treatment 仍有症狀
    \item 人類肺臟由 vagus nerve 的 pulmonary branches 廣泛支配
    \item 在 COPD 中，autonomic nervous system 失調
    \item Airway smooth muscle tone（受 vagal control）導致過度支氣管收縮及 hyperinflation
\end{itemize}

\subsection{Targeted Lung Denervation (TLD)}

\begin{itemize}
    \item 單次門診手術，在全身麻醉下進行
    \item 使用 D'Nerva system 在 mainstem bronchi 交界處進行 circumferential radiofrequency (RF) ablation
    \item 目的：永久性阻斷 parasympathetic nerve signaling，減少 neural hyperactivity
    \item 先前研究顯示可減少 COPD exacerbations 並穩定肺功能
\end{itemize}

%==============================================================
\section{研究方法}
%==============================================================

\subsection{研究設計}

\begin{longtable}{p{4cm}p{10cm}}
\toprule
\textbf{項目} & \textbf{內容} \\
\midrule
\endfirsthead

\toprule
\textbf{項目} & \textbf{內容} \\
\midrule
\endhead

\bottomrule
\endfoot

\textbf{設計} & Prospective, multicenter, 1:1 randomized, sham-controlled, double-blind \\
\midrule
\textbf{Randomization} & 388 patients（TLD: 198, Sham: 190）\\
\midrule
\textbf{研究地點} & 32 sites（美國 19 sites，歐洲及英國 12 sites）\\
\midrule
\textbf{研究期間} & September 2019 - August 2023 \\
\midrule
\textbf{追蹤期間} & 12 months \\
\bottomrule
\end{longtable}

\subsection{納入條件}

\begin{itemize}
    \item 年齡 $\geq$40 歲
    \item 吸菸史 $\geq$10 pack-years
    \item COPD 診斷，moderate to severe airflow obstruction：
    \begin{itemize}
        \item Postbronchodilator FEV$_1$ $\geq$25\% 但 $\leq$80\% predicted
        \item FEV$_1$/FVC ratio <70\%
    \end{itemize}
    \item 前 12 個月有 $\geq$2 moderate 或 $\geq$1 severe exacerbation（GOLD class E）
    \item CAT score $\geq$10
    \item 接受 optimal medical treatment（至少 LABA/LAMA dual therapy $\pm$ ICS）
\end{itemize}

\subsection{排除條件}

\begin{itemize}
    \item Whole-lung emphysema >50\%（density mask threshold <-950 HU）
    \item Symptomatic gastric motility disorder（Gastroparesis Cardinal Symptom Index >18）
    \item Dominant non-COPD lung disease（如 cystic fibrosis、asthma）
\end{itemize}

\subsection{介入方式}

\begin{longtable}{p{4cm}p{10cm}}
\toprule
\textbf{組別} & \textbf{處置} \\
\midrule
\endfirsthead

\toprule
\textbf{組別} & \textbf{處置} \\
\midrule
\endhead

\bottomrule
\endfoot

\textbf{TLD Group} & Active treatment：fluoroscopy + RF energy delivery + 持續 optimal medical treatment \\
\midrule
\textbf{Sham Group} & D'Nerva catheter 置入 + balloon inflation，但無 fluoroscopy 或 RF energy + 持續 optimal medical treatment \\
\bottomrule
\end{longtable}

%==============================================================
\section{研究終點}
%==============================================================

\subsection{Primary Endpoint}

Time to first moderate or severe COPD exacerbation through 12 months

\note{
\textbf{COPD Exacerbation 定義：}
\begin{itemize}
    \item 呼吸症狀（cough、sputum、wheezing、dyspnea、chest tightness）新發或惡化
    \item 至少一種症狀持續 $\geq$3 天
    \item 需要 antibiotics 及/或 systemic steroids 治療（moderate）
    \item 需要住院（severe）
\end{itemize}
}

\subsection{Secondary Endpoints}

\begin{itemize}
    \item Time to first severe COPD exacerbation (0-12 mo)
    \item 12-month change in SGRQ-C score
    \item 12-month change in FVC, FEV$_1$, RV
    \item 12-month change in SF-36 score
    \item 12-month TDI score
    \item Proportion achieving MCID in CAT at 12 months
\end{itemize}

%==============================================================
\section{主要結果}
%==============================================================

\subsection{Baseline Characteristics}

兩組基線特徵良好匹配：
\begin{itemize}
    \item 平均年齡：67.8 vs 67.7 歲
    \item 女性：56.1\% vs 54.2\%
    \item 平均 FEV$_1$：41.6\% vs 40.5\% predicted
    \item GOLD III 佔多數：64.1\% vs 70.0\%
    \item 前一年 exacerbations：2.9 vs 3.0 次
    \item 使用 Triple therapy (LABA+LAMA+ICS)：91.9\% vs 93.2\%
\end{itemize}

\subsection{Primary Endpoint 結果}

\note{
\textbf{Primary Endpoint 未達成：}

\textbf{Hazard Ratio：}1.27 (95\% CI, 0.988-1.628)

\textbf{兩組無統計顯著差異}
}

\begin{itemize}
    \item Cox proportional hazards model 的 proportional hazards assumption 無效 (P=0.013)
    \item 使用 prespecified 3-12 month window 重新計算：HR 0.793 (95\% CI, 0.534-1.177)，仍無顯著差異
\end{itemize}

\subsection{Secondary Endpoints 結果}

\subsubsection{Dyspnea}

\begin{longtable}{p{5cm}p{4cm}p{4cm}p{1.5cm}}
\toprule
\textbf{指標} & \textbf{TLD} & \textbf{Sham} & \textbf{P} \\
\midrule
\endfirsthead

\toprule
\textbf{指標} & \textbf{TLD} & \textbf{Sham} & \textbf{P} \\
\midrule
\endhead

\bottomrule
\endfoot

TDI $\geq$1-point improvement & \textbf{35.4\%} & 24.1\% & \textbf{0.021} \\
\midrule
mMRC change from baseline & -0.25$\pm$0.96 & -0.02$\pm$0.93 & \textbf{0.006} \\
\midrule
mMRC $\geq$1-point decrease & 35.4\% & 25.8\% & 0.051 \\
\bottomrule
\end{longtable}

\subsubsection{Lung Function}

\begin{longtable}{p{3cm}p{3cm}p{3cm}p{3cm}p{2cm}}
\toprule
\textbf{參數} & \textbf{TLD Change} & \textbf{P vs baseline} & \textbf{Sham Change} & \textbf{P vs baseline} \\
\midrule
\endfirsthead

\toprule
\textbf{參數} & \textbf{TLD Change} & \textbf{P vs baseline} & \textbf{Sham Change} & \textbf{P vs baseline} \\
\midrule
\endhead

\bottomrule
\endfoot

FVC (L) & -0.01$\pm$0.41 & 0.720 & -0.07$\pm$0.45 & 0.039 \\
\midrule
FEV$_1$ (L) & -0.01$\pm$0.15 & 0.281 & \textbf{-0.05$\pm$0.17} & \textbf{<0.001} \\
\midrule
RV (L) & -0.02$\pm$1.03 & 0.791 & +0.16$\pm$0.88 & 0.017 \\
\bottomrule
\end{longtable}

\note{
\textbf{重要發現：}
\begin{itemize}
    \item \textbf{TLD 組：}FEV$_1$ 12 個月內無下降（P=0.281 vs baseline）
    \item \textbf{Sham 組：}FEV$_1$ 顯著下降（P<0.001 vs baseline）
    \item COPD 病人預期 FEV$_1$ 每年下降約 50 mL
\end{itemize}
}

\subsubsection{Quality of Life}

\begin{itemize}
    \item SGRQ-C：兩組均較 baseline 改善，但組間無顯著差異
    \item CAT score：TLD 組改善 -2.2$\pm$6.8，Sham 組改善 -1.2$\pm$7.3（P=0.104）
\end{itemize}

%==============================================================
\section{Post Hoc 分析：Responder Phenotype}
%==============================================================

\note{
\textbf{關鍵發現：}Post hoc 分析識別出 TLD 治療的 responder profile：

\textbf{Airway-Predominant Phenotype：}
\begin{itemize}
    \item Lung hyperinflation（高 RV\% predicted）
    \item 無顯著 emphysema（低 \% emphysema score）
\end{itemize}
}

\subsection{分析結果}

\begin{itemize}
    \item Baseline emphysema 程度與 12-month FEV$_1$ 改善呈負相關（Pearson r = -0.21 for TLD）
    \item 在 baseline hyperinflation 較高且 emphysema 較低的病人中：
    \begin{itemize}
        \item FEV$_1$ 改善幅度較大
        \item COPD exacerbation 減少幅度較大
    \end{itemize}
    \item 這些發現與 patient sex、age、blood eosinophilia、study site/region 無關
\end{itemize}

%==============================================================
\section{安全性}
%==============================================================

\subsection{Serious Adverse Events (0-12 months)}

\begin{longtable}{p{6cm}p{4cm}p{4cm}}
\toprule
\textbf{項目} & \textbf{TLD} & \textbf{Sham} \\
\midrule
\endfirsthead

\toprule
\textbf{項目} & \textbf{TLD} & \textbf{Sham} \\
\midrule
\endhead

\bottomrule
\endfoot

Overall SAE rate & 37.4\% & 34.2\% \\
\midrule
COPD exacerbation & 21.7\% & 15.8\% \\
\midrule
Respiratory failure & 2.5\% & 3.2\% \\
\midrule
GI disorders & 6.6\% & 3.7\% \\
\midrule
All-cause mortality & 4.5\% & 4.2\% \\
\midrule
Respiratory-related mortality & 2.0\% & 3.7\% \\
\bottomrule
\end{longtable}

\subsection{Periprocedural SAEs (0-3 months)}

\begin{itemize}
    \item TLD 組 SAE rate 較高：\textbf{21.7\% vs 10.5\%}
    \item 主要由 respiratory 及 GI SAEs 驅動
    \item GI SAEs：TLD 5.1\% vs Sham 1.1\%（GI events 是 TLD 的已知併發症）
    \item 1 例 bronchoesophageal fistula（0.5\%），導入 esophageal cooling device 後無再發生
\end{itemize}

%==============================================================
\section{Discussion}
%==============================================================

\subsection{Primary Endpoint 未達成的原因}

\begin{itemize}
    \item TLD 組在術後前 3 個月有較高 exacerbation rate
    \item 與其他 bronchoscopic intervention 試驗觀察一致
    \item AIRFLOW-3 有 42\% 病人在前一年有 hospitalization，高於 AIRFLOW-2（24\%）
    \item Prior-year severe exacerbation history 與 periprocedural respiratory SAE 相關
\end{itemize}

\subsection{正面發現}

\begin{itemize}
    \item \textbf{肺功能保留：}TLD 組 FEV$_1$ 無下降，consistent with previous TLD trials
    \item \textbf{Dyspnea 改善：}TDI 及 mMRC 均顯示改善
    \item \textbf{Responder phenotype 識別：}Airway-predominant phenotype（hyperinflation without significant emphysema）
\end{itemize}

%==============================================================
\section{Take-Home Message}
%==============================================================

\begin{enumerate}
    \item \textcolor{red}{\textbf{Primary Endpoint 未達成：}}AIRFLOW-3 未達到主要終點（減少 time to first moderate/severe COPD exacerbation）
    \item \textcolor{red}{\textbf{肺功能保留：}}TLD 組展現肺功能保留，Sham 組則有預期的 FEV$_1$ 下降
    \item \textcolor{red}{\textbf{Dyspnea 改善：}}TLD 組有較高比例達到 dyspnea MCID（TDI 35.4\% vs 24.1\%）
    \item \textcolor{red}{\textbf{Responder Phenotype：}}Post hoc 分析識別出 responder profile：hyperinflation without significant emphysema（airway-predominant phenotype）
    \item \textcolor{red}{\textbf{未來方向：}}正在設計 prospective RCT 以確認在 responder phenotype 病人中的效益
\end{enumerate}

%==============================================================
\section{縮寫對照表}
%==============================================================

\begin{longtable}{p{3cm}p{11cm}}
\toprule
\textbf{縮寫} & \textbf{全名} \\
\midrule
\endfirsthead

\toprule
\textbf{縮寫} & \textbf{全名} \\
\midrule
\endhead

\bottomrule
\endfoot

TLD & Targeted Lung Denervation (標靶肺去神經術) \\
COPD & Chronic Obstructive Pulmonary Disease \\
RF & Radiofrequency \\
GOLD & Global Initiative for Chronic Obstructive Lung Disease \\
CAT & COPD Assessment Test \\
LABA & Long-Acting Beta-Agonist \\
LAMA & Long-Acting Muscarinic Antagonist \\
ICS & Inhaled Corticosteroid \\
FEV$_1$ & Forced Expiratory Volume in 1 second \\
FVC & Forced Vital Capacity \\
RV & Residual Volume \\
TDI & Transitional Dyspnea Index \\
mMRC & Modified Medical Research Council \\
SGRQ-C & St. George's Respiratory Questionnaire for COPD \\
MCID & Minimal Clinically Important Difference \\
SAE & Serious Adverse Event \\
GI & Gastrointestinal \\
ITT & Intention-To-Treat \\
HR & Hazard Ratio \\
CI & Confidence Interval \\
\bottomrule
\end{longtable}

%==============================================================
\section{參考文獻}
%==============================================================

\begin{enumerate}
    \item Shah PL, Slebos DJ, Sue R, et al. Randomized sham-controlled trial of targeted lung denervation in patients with chronic obstructive pulmonary disease (AIRFLOW-3). \href{https://doi.org/10.1164/rccm.202502-0404OC}{\textit{Am J Respir Crit Care Med}. 2025;211:2318-2329.}

    \item Slebos DJ, Shah PL, Herth FJF, et al. Safety and adverse events after targeted lung denervation for symptomatic moderate to severe chronic obstructive pulmonary disease (AIRFLOW). \href{https://doi.org/10.1164/rccm.201903-0624OC}{\textit{Am J Respir Crit Care Med}. 2019;200:1477-1486.}

    \item Valipour A, Shah PL, Herth FJ, et al. Two-year outcomes for the double-blind, randomized, sham-controlled study of targeted lung denervation in patients with moderate to severe COPD: AIRFLOW-2. \href{https://doi.org/10.2147/COPD.S267409}{\textit{Int J Chron Obstruct Pulmon Dis}. 2020;15:2807-2816.}

    \item Lipson DA, Barnhart F, Brealey N, et al. Once-daily single-inhaler triple versus dual therapy in patients with COPD. \href{https://doi.org/10.1056/NEJMoa1713901}{\textit{N Engl J Med}. 2018;378:1671-1680.}

    \item Singh D, Agusti A, Anzueto A, et al. Global strategy for the diagnosis, management, and prevention of chronic obstructive lung disease: the GOLD science committee report 2019. \href{https://doi.org/10.1183/13993003.00164-2019}{\textit{Eur Respir J}. 2019;53:1900164.}
\end{enumerate}

\end{document}
